% paper101_appendix_scanner.tex
% Supplementary Appendix: CRMLint Scanner Classification of Scholze's
% Berkovich Motives and Motivic Geometrization
% Paper 101 of the Constructive Reverse Mathematics Series
%
% Paul Chun-Kit Lee, NYU
% March 2026

\documentclass[11pt]{article}

\usepackage[margin=1in]{geometry}
\usepackage{booktabs}
\usepackage{amsmath,amssymb}
\usepackage{hyperref}
\usepackage{xcolor}
\usepackage{longtable}
\usepackage{array}
\usepackage[T1]{fontenc}

\newcommand{\BISH}{\textsf{BISH}}
\newcommand{\WKL}{\textsf{WKL}}
\newcommand{\WLPO}{\textsf{WLPO}}
\newcommand{\LPO}{\textsf{LPO}}
\newcommand{\CLASS}{\textsf{CLASS}}

\title{Supplementary Appendix: CRMLint Scanner Output\\
for Berkovich Motives and Motivic Geometrization}
\author{Paul Chun-Kit Lee\\
\small New York University}
\date{March 2026}

\begin{document}
\maketitle

\appendix
\section{CRMLint Scanner Methodology}
\label{sec:methodology}

The CRMLint scanner (Paper~76 of this series; DOI:~10.5281/zenodo.18779362) is an
automated static analysis tool that assigns a preliminary CRM level to each logical
step in a mathematical text.  The scanner used here operates in \emph{keyword mode}
on raw \LaTeX\ source, employing three detection layers:

\begin{enumerate}
\item \textbf{Keyword markers ($\sim$80 patterns).}
  Each pattern maps a mathematical term to its minimal CRM level.
  Examples: ``motivic sheaf'' $\to$ \BISH, ``normed completion'' $\to$ \WLPO,
  ``perfectoid'' $\to$ \WKL, ``normed filtration'' $\to$ \LPO.
  Patterns are case-insensitive substrings applied to each line of source.

\item \textbf{Citation patterns ($\sim$120 entries).}
  Bibliographic references are matched against a curated database of known
  CRM-relevant works.  For example, a citation of Scholze's
  \emph{Lectures on Analytic Geometry} triggers a \WKL\ alert (Tychonoff
  gates require Weak K\"onig's Lemma for compactness of profinite
  inverse limits).

\item \textbf{Anti-marker system.}
  Certain patterns \emph{suppress} false positives.
  For instance, ``algebraic $K$-theory'' in a purely algebraic context
  does not trigger an analytic alert.
  The scanner reports suppressed hits separately (7 in BerkovichMotives,
  3 in MotivicGeometrization).
\end{enumerate}

Scanner output is a JSON file listing each active hit with its line number,
matched pattern, detected CRM level, and the enclosing section title.
Human classification then assigns each hit to one of four categories:

\begin{description}
\item[Category~A (motivational):] Alerts in introductions or preambles.
  These mention analytic infrastructure for context but do not contribute
  to the logical content of any proof.

\item[Category~B (parasitic):] Alerts where an analytic formalism
  (Banach ring completion, seminorm filtration) \emph{carries} algebraic
  content without contributing logical substance.  The WLPO/LPO cost is
  an artifact of the presentation, not the theorem.

\item[Category~\BISH\ (algebraic core):] Alerts confirming that the
  flagged content (motivic sheaves, Tate twists, six-functor formalism)
  is genuinely constructive, operating within Bishop-style mathematics.

\item[Category~C (structural WKL):] Alerts where Weak K\"onig's Lemma
  is \emph{structurally necessary}---inverse limits of profinite objects,
  perfectoid tilting, condensed mathematics infrastructure.
  These represent the irreducible non-constructive content.
\end{description}

\section{Alert Distribution}
\label{sec:distribution}

\subsection{Distribution by Category}

Table~\ref{tab:category} summarizes the 208 active alerts by category
and source document.

\begin{table}[ht]
\centering
\caption{Alert distribution by category.}
\label{tab:category}
\begin{tabular}{@{}lccc@{}}
\toprule
Category & Berkovich Motives & Motivic Geometrization & Total \\
\midrule
A (motivational)    & 11  & 15  & 26  \\
B (parasitic)       & 45  & 18  & 63  \\
\BISH\ (algebraic)  & 43  & 20  & 63  \\
C (structural \WKL) & 38  & 18  & 56  \\
\midrule
Total               & 137 & 71  & 208 \\
\bottomrule
\end{tabular}
\end{table}

The combined parasitic-plus-motivational count ($\text{A} + \text{B} = 89$,
or 43\% of alerts) represents analytic overhead that does not contribute
to the logical structure of any theorem.  The \BISH\ alerts (30\%) confirm
the algebraic core.  The remaining 56 Category~C alerts (27\%) constitute
the genuine \WKL\ dependency.

\subsection{Distribution by Section Type and CRM Level}

Table~\ref{tab:section-level} cross-tabulates section type
(motivation/definition/proof) against CRM level.

\begin{table}[ht]
\centering
\caption{Alert distribution by section type and CRM level.
  \CLASS\ column included to confirm the zero-alert finding.}
\label{tab:section-level}
\begin{tabular}{@{}lcccccc@{}}
\toprule
Section type & \CLASS & \LPO & \WLPO & \WKL & \BISH & Total \\
\midrule
Motivation  & 0 & 0  & 11 & 15 & 21 & 47  \\
Definition  & 0 & 4  & 17 &  7 &  0 & 28  \\
Proof       & 0 & 4  & 38 & 49 & 42 & 133 \\
\midrule
Total       & 0 & 8  & 66 & 71 & 63 & 208 \\
\bottomrule
\end{tabular}
\end{table}

\paragraph{Key observation.}
The \CLASS\ column is uniformly zero.  Across both documents---a combined
93~pages, 24~theorems, 78~lemmas, 90~proofs, and 87~bibliographic
references---the scanner found no alert at the \CLASS\ level.  No
invocation of the Axiom of Choice, Zorn's Lemma, well-ordering, or
any equivalent principle was detected.  This is the primary evidence
supporting the verdict $\mathrm{CRM}(\text{Scholze's motivic programme})
= \WKL$.

\paragraph{The 8 LPO alerts.}
All 8 \LPO\ alerts arise from seminorm filtrations indexed by
$\mathbb{R}_{\geq 0}$, where decidable comparison $\|x\| \leq r$
triggers \LPO.  These are uniformly parasitic (Category~B): the filtration
is an analytic bookkeeping device, and the algebraic content of the
theorems it appears in does not require decidable norm comparison.

\section{Representative Classified Alerts}
\label{sec:samples}

Table~\ref{tab:samples} lists 14 representative alerts spanning all four
categories, both source documents, and multiple CRM levels.

\begin{longtable}{@{}p{2.0cm}cp{2.8cm}cp{5.5cm}@{}}
\caption{Representative classified alerts.  Line numbers refer to the
  \LaTeX\ source.}
\label{tab:samples} \\
\toprule
File & Line & Reason & Cat. & Justification \\
\midrule
\endfirsthead
\toprule
File & Line & Reason & Cat. & Justification \\
\midrule
\endhead
\midrule
\multicolumn{5}{r}{\emph{continued on next page}} \\
\bottomrule
\endfoot
\bottomrule
\endlastfoot
%
% Category A: motivational
%
Berkovich\-Motives & 60
  & [\WKL] Lectures on Analytic Geometry (citation)
  & A
  & Introductory citation establishing background for condensed mathematics. \\
%
Berkovich\-Motives & 194
  & [\WLPO] normed completion
  & A
  & Describing classical Berkovich approach in introduction; not in any proof. \\
%
Motivic\-Geom & 97
  & [\WKL] perfectoid (inverse limit / tilting)
  & A
  & Perfectoid keyword in introduction; motivational context only. \\
\midrule
%
% Category B: parasitic
%
Berkovich\-Motives & 307
  & [\LPO] normed filtration
  & B
  & Seminorm filtration $\mathbb{R}_{\geq 0}$ indexing in Banach ring definition;
    analytic formalism parasitic on algebraic content. \\
%
Berkovich\-Motives & 571
  & [\WLPO] normed completion
  & B
  & Banach ring formalism carried through arc-topology proofs; parasitic. \\
%
Motivic\-Geom & 325
  & [\LPO] normed filtration
  & B
  & Seminorm filtration in Basic results; $\mathbb{R}_{\geq 0}$ indexing
    is parasitic on algebraic content. \\
%
Motivic\-Geom & 498
  & [\WLPO] normed completion
  & B
  & Banach ring formalism in Stacks of $L$-parameters; parasitic. \\
\midrule
%
% Category BISH: algebraic core
%
Berkovich\-Motives & 1172
  & [\BISH] Tate twist (algebraic)
  & \BISH
  & Genuine algebraic marker; the Tate twist is a purely algebraic operation. \\
%
Berkovich\-Motives & 1924
  & [\BISH] six-functor formalism
  & \BISH
  & Genuine algebraic marker in $D_\mathrm{mot}$; six-functor formalism
    is categorical/constructive. \\
%
Motivic\-Geom & 391
  & [\BISH] six-functor formalism
  & \BISH
  & Genuine algebraic marker in $D_\mathrm{mot}$; categorical infrastructure
    is \BISH. \\
\midrule
%
% Category C: structural WKL
%
Berkovich\-Motives & 501
  & [\WKL] inverse limit of finite objects
  & C
  & Genuine profinite/inverse limit construction in the nonarchimedean
    unit disc definition. \\
%
Berkovich\-Motives & 665
  & [\WKL] perfectoid (tilting)
  & C
  & Structural: perfectoid inverse limit/tilting in the arc-topology proofs. \\
%
Berkovich\-Motives & 1607
  & [\WKL] Lectures on Analytic Geometry (citation)
  & C
  & Structural dependency on condensed infrastructure in $K$-theory
    proof section. \\
%
Motivic\-Geom & 436
  & [\WKL] perfectoid (tilting)
  & C
  & Structural: perfectoid inverse limit in $D_\mathrm{mot}$ proof section. \\
\end{longtable}

\section{Zero-CLASS Confirmation}
\label{sec:zero-class}

The central finding of this scanner analysis is the \emph{absence} of
\CLASS-level alerts.  Specifically:

\begin{itemize}
\item The keyword database includes markers for the Axiom of Choice
  (``Axiom of Choice,'' ``Zorn's Lemma,'' ``well-ordering,''
  ``transfinite induction,'' ``uncountable choice''), Baire category
  (``Baire category theorem,'' ``generic point argument''),
  and other \CLASS\ indicators.
  None triggered.

\item The citation database includes 17 works known to require full
  \CLASS\ (e.g., general Hahn--Banach, uncountable Zorn applications).
  None were cited.

\item The anti-marker system did not suppress any \CLASS-level hits.
  The 10 total suppressions (7 + 3) were all at \BISH\ or \WKL\ levels.
\end{itemize}

Combined with the structural observation that Scholze's programme
replaces point-set topology (which requires \CLASS\ for Tychonoff and
Stone--\v{C}ech) with condensed mathematics (which requires only \WKL\
for profinite compactness), the zero-\CLASS\ finding is not accidental
but architectural.

\section{Reproduction Instructions}
\label{sec:reproduction}

The scanner and classification can be reproduced as follows.
All files are in the directory
\texttt{paper~101/discussion\_paper101\_appendix\_prelim/}.

\begin{enumerate}
\item \textbf{Run the CRMLint scanner on each source file.}

\begin{verbatim}
python3 crmlint_scanner.py BerkovichMotives.tex \
    -o berkovich_scan.json --format json

python3 crmlint_scanner.py MotivicGeometrization.tex \
    -o motivic_geom_scan.json --format json
\end{verbatim}

Expected output: 137 active hits for BerkovichMotives
(\BISH=43, \WKL=47, \WLPO=40, \LPO=7)
and 71 active hits for MotivicGeometrization
(\BISH=20, \WKL=24, \WLPO=26, \LPO=1).

\item \textbf{Classify alerts into categories.}

The classification follows the rules described in
Section~\ref{sec:methodology}.  The output is stored in
\texttt{scanner\_classification.csv} with columns:
\texttt{file}, \texttt{line}, \texttt{level}, \texttt{reason},
\texttt{location}, \texttt{category}, \texttt{section\_type},
\texttt{justification}.

\item \textbf{Verify totals.}

The CSV should contain exactly 208 data rows:
26 Category~A, 63 Category~B, 63 Category~\BISH,
56 Category~C.  The \CLASS\ column of
Table~\ref{tab:section-level} should be uniformly zero.
\end{enumerate}

\bigskip

\noindent\textbf{Scanner version.}
CRMLint scanner v1.0 (Paper~76, DOI:~10.5281/zenodo.18779362).
Keyword database: 80 patterns, 120 citation entries.
Python~3.11+, no external dependencies beyond the standard library.

\end{document}
